\chapter*{Glossary and Abbreviations}
\addcontentsline{toc}{chapter}{Glossary and Abbreviations}

\begin{description}
    \item[API] Application Programming Interface. A set of rules that allows different software entities to communicate.
    \item[DTO] Data Transfer Object. An object that carries data between processes to reduce the number of method calls and network overhead.
    \item[EF Core] Entity Framework Core. An open-source, lightweight, and extensible object-relational mapper (O/RM) for the .NET ecosystem.
    \item[JWT] JSON Web Token. An open standard (RFC 7519) that defines a compact and self-contained way for securely transmitting information between parties as a JSON object.
    \item[LLM] Large Language Model. A deep learning algorithm, utilizing transformer architecture, that can recognize, summarize, translate, predict, and generate human-like text.
    \item[POI] Point of Interest. A specific geographical location (e.g., restaurant, museum, or landmark) that a user may find useful or interesting.
    \item[RAG] Retrieval-Augmented Generation. An AI architectural framework that retrieves data from external knowledge bases and injects it into an LLM's context window to ground the generation and prevent hallucinations.
    \item[SERP] Search Engine Results Page. The page displayed by a search engine in response to a query, previously utilized for targeted image fetching.
    \item[UI/UX] User Interface / User Experience. The graphical layout and the holistic interactive experience provided to the end-user by the application.
    \item[UUID] Universally Unique Identifier. A 128-bit label used for information in computer systems, heavily utilized as Primary Keys in the Kemora database.
\end{description}
